\section{Rebalancing by rotation}
\label{sec:rotations}

\Cref{sec:height} shows what the balance condition is worth once it holds. This
section gives the operation that restores it after an update has broken it. The
operation is called a rotation, it rearranges three links, and it is the only
repair an AVL tree ever performs. Everything in \Cref{sec:insertion} and
\Cref{sec:deletion} is a question of where to apply it and how often.

\subsection{The operation}

A right rotation at a node $z$ takes the left child $y$ of $z$ and makes it the
new root of the subtree, with $z$ descending to become the right child of $y$.
The subtree that $y$ previously held on its right becomes the new left subtree of
$z$, which is the only place it can go and still satisfy \Cref{def:bst}. A left
rotation is the mirror image. \Cref{fig:rotation} shows the right rotation on the
smallest tree that needs one.

\begin{figure}[t]
  \centering
  % A single right rotation, before and after, with the three links that change
% drawn heavy and everything else left alone.
%
% The point of the lower rail is that the left-to-right reading of the keys is
% identical on both sides. A rotation moves nodes vertically; it never moves
% them past one another.
\begin{tikzpicture}[x = 1mm, y = 1mm]

  % ------------------------------------------------------------ before
  \begin{scope}
    \node[tmark]  (z) at (  8,   0) {30};
    \node[tnode]  (y) at ( -6, -14) {20};
    \node[tnode]  (x) at (-20, -28) {10};
    \draw[tedge, line width = 1.1pt] (z) -- (y);
    \draw[tedge, line width = 1.1pt] (y) -- (x);
    \node[tlabel, right = 2mm of z] {$\bfac = +2$};
    \node[tannot] at (-6, -42) {before};
  \end{scope}

  % ------------------------------------------------------------- arrow
  \draw[-{Stealth[length=3mm]}, draw = rmid, line width = 0.7pt]
    (34, -14) -- (50, -14)
    node[midway, above, tlabel] {rotate right at 30};

  % ------------------------------------------------------------- after
  \begin{scope}[xshift = 74mm]
    \node[tnode] (ny) at (  0,  -7) {20};
    \node[tnode] (nx) at (-14, -21) {10};
    \node[tnode] (nz) at ( 14, -21) {30};
    \draw[tedge, line width = 1.1pt] (ny) -- (nx);
    \draw[tedge, line width = 1.1pt] (ny) -- (nz);
    \node[tlabel, right = 2mm of ny] {$\bfac = 0$};
    \node[tannot] at (0, -42) {after};
  \end{scope}

  % ----------------------------------------- the reading order, both sides
  \draw[tfaintedge] (-26, -52) -- (14, -52);
  \foreach \k [count = \i from 0] in {10, 20, 30}
    \node[tlabel] at (-26 + \i * 20, -56) {\k};

  \draw[tfaintedge] (60, -52) -- (100, -52);
  \foreach \k [count = \i from 0] in {10, 20, 30}
    \node[tlabel] at (60 + \i * 20, -56) {\k};

  \node[tannot] at (37, -64) {the in-order sequence is the same on both sides};

\end{tikzpicture}

  \caption{A right rotation at the node $30$. Three links change, the node $20$
  rises, and the node $30$ descends. The rail beneath each tree gives the
  in-order reading of the keys, which the rotation leaves alone.}
  \label{fig:rotation}
\end{figure}

\begin{algorithm}[t]
\caption{Rotations. Each one rewrites three links and recomputes two heights.}
\label{alg:rotate}
\SetKwFunction{RotR}{RotateRight}
\SetKwFunction{RotL}{RotateLeft}
\SetKwFunction{Upd}{UpdateHeight}
\Fn{\RotR{$z$}}{
  $y \gets z.\mathit{left}$\;
  $z.\mathit{left} \gets y.\mathit{right}$ \Comment*[r]{the middle subtree changes parent}
  $y.\mathit{right} \gets z$\;
  \Upd{$z$}\; \Comment*[r]{$z$ is now the lower of the two}
  \Upd{$y$}\;
  \KwRet $y$ \Comment*[r]{the new root of this subtree}
}
\BlankLine
\Fn{\RotL{$z$}}{
  $y \gets z.\mathit{right}$\;
  $z.\mathit{right} \gets y.\mathit{left}$\;
  $y.\mathit{left} \gets z$\;
  \Upd{$z$}\;
  \Upd{$y$}\;
  \KwRet $y$\;
}
\end{algorithm}

Two properties make the rotation usable as a repair. It must not destroy the
search tree property, and it must be cheap.

\begin{lemma}[Rotation preserves order]
\label{lem:inorder}
Let $T'$ be the tree obtained from a binary search tree $T$ by a single rotation.
Then $T'$ is a binary search tree, and its in-order key sequence equals that of
$T$.
\end{lemma}

\begin{proof}
Consider a right rotation at $z$ with left child $y$, and name the three subtrees
$A = L(y)$, $B = R(y)$ and $C = R(z)$. Reading $T$ in order visits $A$, then $y$,
then $B$, then $z$, then $C$, because $y$ precedes $z$ and $B$ lies between them.
After the rotation, $y$ is the root, $A$ is its left subtree, and its right
subtree is the node $z$ carrying $B$ on the left and $C$ on the right. Reading
$T'$ in order visits $A$, then $y$, then $B$, then $z$, then $C$, which is the
same sequence. A binary tree whose in-order traversal is sorted satisfies
\Cref{def:bst}, and the sequence was sorted in $T$, so it is sorted in $T'$. The
left rotation is the mirror argument.
\end{proof}

\Cref{lem:inorder} is the reason a rotation is safe to apply anywhere. It moves
nodes vertically and never past one another horizontally, so no key changes its
rank.

\begin{lemma}[Rotation cost]
\label{lem:rotcost}
A rotation runs in $O(1)$ time, independent of the size of the subtrees involved.
\end{lemma}

\begin{proof}
\Cref{alg:rotate} performs three link assignments and two height updates, each of
which reads the heights of two children and takes a maximum. The subtrees $A$,
$B$ and $C$ are moved by reassigning a pointer to their roots, and none of their
interior nodes is visited. The work is therefore bounded by a constant.
\end{proof}

\subsection{The four cases}

A single rotation is not always enough. Which repair applies depends on the shape
of the path from the unbalanced node down toward the subtree that grew, and there
are four shapes. Let $z$ be the lowest unbalanced node, let $y$ be the child of
$z$ on the taller side, and let $x$ be the child of $y$ on its taller side. The
case is named for the two steps taken from $z$ down to $x$.

\begin{table}[t]
  \centering
  \caption{The four imbalance cases and the repair each one needs. The first two
  lean consistently in one direction; the last two change direction on the second
  step, and the first rotation of the pair straightens them into one of the first
  two.}
  \label{tab:cases}
  \begin{tabular}{@{}llll@{}}
    \toprule
    Case & Path from $z$ to $x$ & Repair & Single rotations \\
    \midrule
    \case{ll} & left, then left   & right rotation at $z$                  & 1 \\
    \case{rr} & right, then right & left rotation at $z$                   & 1 \\
    \case{lr} & left, then right  & left at $y$, then right at $z$         & 2 \\
    \case{rl} & right, then left  & right at $y$, then left at $z$         & 2 \\
    \bottomrule
  \end{tabular}
\end{table}

\Cref{fig:fourcases} draws all four on three keys. The \case{ll} and \case{rr}
cases lean the same way twice, and one rotation at $z$ lifts the middle key to
the top. The \case{lr} and \case{rl} cases zig-zag, and a single rotation at $z$
would only move the imbalance to the other side. The first rotation of the pair
is applied at $y$ rather than $z$, and it converts the zig-zag into a straight
lean, after which the second rotation finishes the repair exactly as in the
straight case. All four end at the same tree.

\begin{figure}[t]
  \centering
  % The four ways an insertion can break a node, and the shape each one leaves
% behind. The heavy node is the lowest node whose balance factor has reached
% two; the arrow marks the node the first rotation turns about.
%
% Read the pairs across: LL and RR are mirror images, and so are LR and RL.
% Read them down: the second row needs a preparatory turn that the first does
% not, because its three nodes zig-zag instead of leaning one way.
\begin{tikzpicture}[x = 1mm, y = 1mm]

  % --------------------------------------------------------------- LL
  \begin{scope}
    \node[tmark] (c) at ( 10,   0) {30};
    \node[tnode] (b) at ( -2, -12) {20};
    \node[tnode] (a) at (-14, -24) {10};
    \draw[tedge] (c) -- (b);  \draw[tedge] (b) -- (a);
    \node[tlabel] at (-2, -33) {\case{ll}: one right rotation at 30};
  \end{scope}

  % --------------------------------------------------------------- RR
  \begin{scope}[xshift = 58mm]
    \node[tmark] (C) at (-14,   0) {10};
    \node[tnode] (B) at ( -2, -12) {20};
    \node[tnode] (A) at ( 10, -24) {30};
    \draw[tedge] (C) -- (B);  \draw[tedge] (B) -- (A);
    \node[tlabel] at (-2, -33) {\case{rr}: one left rotation at 10};
  \end{scope}

  % --------------------------------------------------------------- LR
  \begin{scope}[yshift = -46mm]
    \node[tmark] (d) at ( 10,   0) {30};
    \node[tnode] (e) at (-14, -12) {10};
    \node[tnode] (f) at ( -2, -24) {20};
    \draw[tedge] (d) -- (e);  \draw[tedge] (e) -- (f);
    \node[tlabel] at (-2, -33) {\case{lr}: left at 10, then right at 30};
  \end{scope}

  % --------------------------------------------------------------- RL
  \begin{scope}[xshift = 58mm, yshift = -46mm]
    \node[tmark] (D) at (-14,   0) {10};
    \node[tnode] (E) at ( 10, -12) {30};
    \node[tnode] (F) at ( -2, -24) {20};
    \draw[tedge] (D) -- (E);  \draw[tedge] (E) -- (F);
    \node[tlabel] at (-2, -33) {\case{rl}: right at 30, then left at 10};
  \end{scope}

  % ----------------------------------------------- what all four settle into
  \begin{scope}[xshift = 112mm, yshift = -23mm]
    \node[tnode] (t) at (  0,   0) {20};
    \node[tnode] (l) at (-12, -13) {10};
    \node[tnode] (r) at ( 12, -13) {30};
    \draw[tedge] (t) -- (l);  \draw[tedge] (t) -- (r);
    \node[tannot, text width = 30mm, align = center] at (0, -28)
      {all four settle into this};
  \end{scope}

\end{tikzpicture}

  \caption{The four imbalance cases on three keys. The heavy node is the lowest
  unbalanced one. Whichever order the keys arrive in, the repair produces the
  same balanced tree.}
  \label{fig:fourcases}
\end{figure}

\Cref{alg:rebalance} selects among the cases. It inspects the balance factor of
$z$ to find the taller side, then the balance factor of $y$ to find whether the
lean is straight or a zig-zag.

\begin{algorithm}[t]
\caption{Rebalancing one node. The node is returned because a rotation changes
which node roots the subtree, and the caller must relink it.}
\label{alg:rebalance}
\SetKwFunction{Reb}{Rebalance}
\SetKwFunction{RotR}{RotateRight}
\SetKwFunction{RotL}{RotateLeft}
\SetKwFunction{Upd}{UpdateHeight}
\Fn{\Reb{$z$}}{
  \Upd{$z$}\;
  \uIf{$\bfac(z) > 1$}{
    \Comment*[l]{left side is too tall}
    \lIf{$\bfac(z.\mathit{left}) < 0$}{$z.\mathit{left} \gets$ \RotL{$z.\mathit{left}$}}
    \Comment*[l]{\case{lr}: straighten first}
    \KwRet \RotR{$z$} \Comment*[r]{\case{ll}}
  }
  \uElseIf{$\bfac(z) < -1$}{
    \lIf{$\bfac(z.\mathit{right}) > 0$}{$z.\mathit{right} \gets$ \RotR{$z.\mathit{right}$}}
    \Comment*[l]{\case{rl}: straighten first}
    \KwRet \RotL{$z$} \Comment*[r]{\case{rr}}
  }
  \Else{
    \KwRet $z$ \Comment*[r]{already legal, nothing to do}
  }
}
\end{algorithm}

Taken together, \Cref{lem:inorder} and \Cref{lem:rotcost} give the property the
rest of the report depends on. The repair is correct wherever it is applied and
costs a constant, so the cost of an update is decided entirely by how many times
the repair runs. \Cref{sec:insertion} and \Cref{sec:deletion} answer that
question, and they give different answers.
